\section{Our Approach}
\label{sec:impl}

Five configurations grow the same exact CART or C4.5 tree. They differ in how
a node's rows remain sorted, which cuts are scored explicitly, and where and
in what order nodes are grown. \ImplGpu{} and \ImplHybrid{} are two modes of
one CUDA grower, without and with CPU subtree hand-offs.

\subsection{Common Design}
\label{sec:impl-common}

Feature values are stored column-major as 32-bit floats, and labels are mapped
to dense integers. All configurations use the same CART and C4.5 rule layer
and the same double-precision gain calculation, compiled for both CPU and GPU.
Gains within \(10^{-12}\) are ties; the earlier cut and then the lower feature
index win. Optimized searches use conservative filters, but the shared exact
calculation selects the final cut.

\ImplSerialBase{} works with row-identifier lists. The other configurations
represent each feature by presorted 8-byte entries containing its value and a
packed row and class identifier. A node occupies the same interval in every
feature column. A stable partition gives its children sorted intervals, so no
feature is sorted again, following the principle of SLIQ and
SPRINT~\cite{mehta1996sliq,shafer1996sprint}. Concurrent builders first place
nodes in a shared chunked store and finally renumber them into flat preorder
arrays; creation order therefore cannot alter the output. Tests require
identical trees from all configurations. Shared CPU code performs CART or
C4.5 post-processing after growth.

\subsection{\ImplSerialBase}
\label{sec:impl-simple}

\ImplSerialBase{} is a direct sequential implementation and the reference for
the others. At each node, it copies the row identifiers for every feature,
sorts them by value, scores every admissible cut, applies the shared
split-selection rule, and recursively grows two new lists. It therefore
repeats \(O(Fn\log n)\) sorting for a node with \(n\) rows and \(F\) features,
accesses values indirectly through row identifiers, and allocates temporary
vectors throughout growth.

\subsection{\ImplSerialOpt}
\label{sec:impl-serial}

\ImplSerialOpt{} replaces repeated sorting by a stable, single-threaded radix
presort. The first \(n_L\) entries of the winning column form the left child;
the remaining columns are stably partitioned to match, preserving their
order. A feature constant within a node is removed for its whole subtree.

During each sweep, class counts and an inexpensive gain estimate are updated
incrementally. Only candidates that can still beat the incumbent under
conservative rounding and convexity bounds are scored exactly, in cut order;
binary-class specialization and bulk handling of equal values reduce the
remaining work. The winning sweep also supplies the children's class counts.

Columns are allocated once in slots. The serial builder keeps one spare slot:
large ranges move out of place through a one-bit row map, whereas small ranges
use a reused buffer. Large allocations request transparent huge pages, and an
explicit stack replaces recursion. \ImplParallel{} is this same builder with a
thread pool enabled.

\subsection{\ImplParallel}
\label{sec:impl-parallel}

\ImplParallel{} combines two granularities of work. Nodes with at least
32,768 rows sweep features and partition columns in parallel. After a node
with at least 4,096 rows is split, one child becomes a pool task while the
current worker continues with the other. Thus data parallelism serves the few
large nodes near the root, and task parallelism serves the many smaller
subtrees below. Presort and final preorder assembly are parallel as well, and
up to \(\min(F-1,p)\) spare slots let \(p\) threads move columns concurrently.

The CPU path has no socket-aware scheduler or explicit NUMA memory policy:
workers are not pinned and columns are not divided by socket. The pool respects
the process affinity mask, so external \texttt{taskset} or \texttt{numactl}
restrictions remain in force; page placement is otherwise left to Linux
first-touch behavior.

\subsection{\ImplGpu}
\label{sec:impl-gpu}

\ImplGpu{} (\texttt{--cuda --gpu-only}) sets the hand-off threshold to one
row, so every nonempty subtree remains on the GPU. Unlike the depth-first CPU
builders, it processes all active nodes of one level together. Each
node--feature interval is divided into 2,048-entry tiles, one CUDA block per
tile. A level detects constant intervals, forms tile histograms and prefixes,
finds each feature's best cut, and returns these compact results to the host.
The host applies the shared CART or C4.5 choice rule. The device then stably
partitions the columns and simultaneously forms the next level's histograms,
requiring one split-decision round trip per level.
The CPU still stages transfers, chooses splits, and runs shared
post-processing.

The device stores the same entries in one slot per feature plus one spare, and
CUB performs a stable radix presort as columns arrive. On the evaluated A100
and H100, cuts are scored directly in double precision. Devices with slow
double precision may first apply a conservative single-precision filter and
then score every surviving candidate exactly.

Uploads pass through page-locked staging buffers rather than pinning the
training matrix. This is the implementation's only explicit NUMA placement:
on a multi-socket host, the buffers are allocated by a thread bound to the
GPU's PCI NUMA node, and dedicated copy threads are pinned to that node's
CPUs. General training memory and workers are not bound this way.

\subsection{\ImplHybrid}
\label{sec:impl-hybrid}

\ImplHybrid{} retains level-synchronous GPU growth for large nodes but sends
smaller children to the \ImplParallel{} builder
(Fig.~\ref{fig:approach-hybrid}). The threshold is configurable; automatically
it is \(N/8192\), clamped to 1,024--2,048 rows for CART, and \(N/4096\),
clamped to 1,024--4,096 rows for C4.5. These empirical defaults reflect that a
GPU level has substantial fixed overhead, whereas CPU threads can grow small
subtrees independently.

\begin{figure}[t]
\centering
\includegraphics[width=\linewidth]{figs/approach-hybrid.pdf}
\caption{One level of \ImplHybrid{}. Children with at least \(m\) rows form
the next GPU level. Smaller children are gathered into a device ring, copied
through pinned buffers into a host arena, and grown by the CPU pool while the
GPU continues. Both paths contribute to one shared node store.
\ImplGpu{} sets \(m=1\), so no nonempty child is handed off.}
\label{fig:approach-hybrid}
\end{figure}

For a hand-off, the GPU gathers only columns that can still split the child
into a bounded device ring. Copy streams transfer them through pinned buffers
into a pageable host arena, and a host callback starts CPU tasks when their
data arrive. The GPU level loop does not wait for these tasks. The arena is
pre-touched by CPU threads but is not bound to the GPU's NUMA node. CPU and
GPU branches allocate from the same node store and are assembled in preorder
after both finish. Assignment is therefore per node of one tree, not per tree
of a forest, and the CPU grows subtrees rather than merely controlling the
GPU.
